% Introduction to Part IV. Unnumbered text: no \label except the \phantomsection target.
\phantomsection\label{intro:IV}
Parts~I--III treat linear operators. Part~IV constructs the nonlinear input of the variational problem and of the renormalization group: Ba\l aban's symmetric four-dimensional averaging \cite[(0.3)--(0.12)]{BalYM4I}, the averaged axial gauge, the gauge-adapted averaging $\mathcal Q$, the expansion of the Wilson action, and the axial-to-Landau chart. Part~V solves the variational problem in this chart, and the averaging defines the renormalization-group map of Part~VII. From this Part on, constants may depend on $L$ (conventions in Section~\ref{sec:averaging}).

\paragraph{What this Part proves.}
\begin{itemize}
\item \emph{The averaging} (Section~\ref{sec:averaging}). The blocks are centred at the coarse points, and the transport from a block centre to a point of the block is the group average (Lemma~\ref{lem:group-average}) of the transports along the coordinate-ordered contours, one for each order of the directions (Definitions~\ref{def:avg-contour} and~\ref{def:one-step}). The averaging is covariant under the values of gauge transformations at the coarse points and under the lattice symmetries (Lemma~\ref{lem:one-step-basic}), and it preserves small curvature (Proposition~\ref{prop:avg-regular}).
\item \emph{The bridge to the linear theory} (Section~\ref{sec:adapted}). The gauge-adapted averaging $\mathcal Q$ is related to the averaging by the correspondence theorem (Theorem~\ref{thm:correspondence}), and $\mathcal Q_j(e^{\ii\eta A})=e^{\ii\delta_jB_j}$ with $B_j=\bQ_jA+\Cav_j(A)$ (Proposition~\ref{prop:Q-expansion}). The linearized pair $(Q'_k(U_0),\bQ_k(U_0))$ is exactly covariant ($E=0$), local, close to the averages of Part~I, and covariant under $\mathfrak S_k$ (Proposition~\ref{prop:linearized-averages}). So the real results of Parts~II and~III apply to it (Corollary~\ref{cor:parts-apply}).
\item \emph{The action} (Section~\ref{sec:action}). Near a regular $U_0$, $S(e^{\ii\eta A}U_0)=S(U_0)+\ip AJ+\frac12\|D_{U_0}A\|^2+\mathcal N(A)$, with the current $J=J_{U_0}$ and $|\nabla\mathcal N(A)|_\infty\le C(\eps+|A|_\infty)|A|_{C^1}$ (Proposition~\ref{prop:action}).
\item \emph{The chart} (Section~\ref{sec:chart}). Let $U_0$ be $(M,\eps)$-regular with current at most $\eps$, where $\eps$ is small. Every real configuration with small curvature and current whose averages are close to those of $U_0$ has exactly one small representative $e^{\ii\eta A}$ in Ba\l aban's Landau gauge relative to $U_0$; $A$, its covariant gradient, its Hodge Laplacian and the H\"older quotients of its gradient are bounded by $\eps$, the curvature and the current of the configuration, and the difference of the averages (Theorem~\ref{thm:chart}). So the constraint set $\{U:U^{(k)}=V\}$ is parametrized by Lie fields in the Landau gauge (Corollary~\ref{cor:parametrization}).
\end{itemize}
The main theorem is Theorem~\ref{thm:H}, which is Theorem~\ref{thm:chart}: Ba\l aban's axial-to-Landau theorem \cite[Theorem~2]{BalReg} on a single torus.

\paragraph{Imports.} Lemma~\ref{lem:avg} and Theorem~\ref{thm:local-regularity} of Part~I, the perturbation lemmas and Lemmas~\ref{lem:cov-expansion}--\ref{lem:cov-avg} of Part~II, and the regular-background theorems of Part~III with Remark~\ref{rem:general-averaging}. Three passages point ahead: the constants of Lemma~\ref{lem:sup-bounds} are made uniform by Lemma~\ref{lem:uniformity}, whose proof uses only Parts~I--IV, Definition~\ref{def:avg-parameter} and Lemma~\ref{lem:regular-monotone}; the clause of Corollary~\ref{cor:parts-apply} on complex perturbations uses Proposition~\ref{prop:AB}; and one step of the proof of Proposition~\ref{prop:linearized-averages}(f) is written out in the proof of Lemma~\ref{lem:ol-chart-covariance}(a).

\paragraph{Exports.} Propositions~\ref{prop:Q-expansion}, \ref{prop:linearized-averages} and~\ref{prop:action} and Corollary~\ref{cor:parts-apply} are used throughout Parts~V--XI. Theorem~\ref{thm:chart} and Corollary~\ref{cor:parametrization} go to the variational problem (Sections~\ref{sec:var-theorem}--\ref{sec:var-regularity}), the induction (Section~\ref{sec:induction}), the step quantities (Section~\ref{sec:step-quantities}) and the Landau chart of Part~XI (Sections~\ref{sec:ol-composite} and~\ref{sec:ol-chart}). The averaged relative transports $\mathfrak G_W$ (Definition~\ref{def:avg-transport}, Lemma~\ref{lem:avg-transport}) are the linearized averaged contour variables of the gauge fixing in Part~VII. The lattice symmetry of the averaging enters Theorems~\ref{thm:K}(iv), \ref{thm:L}(i) and~\ref{thm:M}(ii). All of Part~IV is on the main line, except Remark~\ref{rem:avg-numerical}, a numerical check that no proof uses.

\paragraph{Details for this Part.} The companion has no chapter on Part~IV; its parameter ledger lists the constants of Part~IV that Parts~VII--XI use.
